\section{Related Work}
\label{sec:related}

\subsection{Generalization Gap in Deep Learning}

The disconnect between benchmark and deployment performance has been extensively documented. \citet{recht2019imagenet} created ImageNetV2 by replicating the original data collection process and found that all tested classifiers experienced 11--14\% accuracy drops. These studies measure gaps through evaluation metrics but do not identify model-internal signals that might predict or explain such gaps.

\subsection{Underspecification and Shortcut Learning}

\citet{damour2020underspecification} formalized underspecification as a fundamental challenge: multiple models can achieve equivalent training performance while exhibiting different behaviors under distribution shift. \citet{wang2025imagenet} showed that ImageNet-trained CNNs learn frequency shortcuts. Our benchmark fingerprint framework generalizes this analysis, asking whether fine-tuning creates detectable signatures regardless of the specific spurious mechanism.

\subsection{Linear Probing for Representation Analysis}

Linear probing has become a standard tool for understanding what information neural network representations encode \citep{alain2017understanding,kornblith2019better}. We apply this methodology to detect benchmark identity from fine-tuned representations.

\subsection{Our Positioning}

Prior work establishes that generalization gaps exist, underspecification causes divergence, and shortcuts can be domain-specific. We contribute a novel representation-level perspective: measuring benchmark fingerprints directly using linear probing. Our negative finding---that BFS does not correlate with gap---opens new research directions.
